% Vietnamese translation for OI Public Library
% Source-File: IOI中国国家候选队论文/2007/day2/4.刘雨辰《对拟阵的初步研究》.ppt
\documentclass[11pt]{article}
\usepackage{oipl-vn}

\newcommand{\Oh}{\mathcal{O}}
\newcommand{\rank}{\operatorname{rank}}

\title{Bản trình chiếu: Nghiên cứu nhập môn về matroid}
\author{Liu Yuchen}
\date{2007}

\begin{document}
\maketitle

\origtitle{Slide companion for an introductory study of matroids}
\sourcepath{IOI China candidate team papers / 2007 / day 2 / Liu Yuchen.ppt}

\section{Phạm vi bản dịch}

Tệp PowerPoint đi cùng bài viết đã được dịch từ PDF. Bản ghi chú này giữ lại
định nghĩa matroid, lý do thuật toán tham lam đúng và các ví dụ chính.

\section{Hệ tập con}

Một hệ tập con là cặp \(M=(S,L)\), trong đó \(S\) là tập hữu hạn và \(L\) là
một họ các tập con hợp lệ. Điều kiện di truyền yêu cầu nếu \(B\in L\) và
\(A\subseteq B\), thì \(A\in L\).

Matroid bổ sung điều kiện trao đổi: nếu \(A,B\in L\) và \(|A|<|B|\), tồn tại
\(x\in B-A\) sao cho \(A\cup\{x\}\in L\). Điều kiện này làm cho mọi tập độc
lập cực đại có cùng kích thước và tạo nền cho thuật toán tham lam.

\section{Tham lam trên matroid}

Nếu mỗi phần tử có trọng số, bài toán chọn tập độc lập có tổng trọng số lớn
nhất được giải bằng cách xét phần tử theo trọng số giảm dần và thêm phần tử
nếu tập vẫn độc lập. Tính trao đổi bảo đảm lựa chọn cục bộ không làm mất khả
năng đạt nghiệm tối ưu toàn cục.

\section{Ví dụ}

Matroid đồ thị lấy \(S\) là tập cạnh và \(L\) là các tập cạnh không tạo chu
trình. Thuật toán tham lam tương ứng chính là Kruskal cho cây khung nhỏ nhất
hoặc lớn nhất tùy cách đặt trọng số. Bài viết cũng nhắc ba lô phân số, lập
lịch và matroid tuyến tính như các trường hợp minh họa cho cùng một cấu trúc
trừu tượng.

\end{document}
